% T2.5 / T2.6 / T2.7 — replacement for Section III "Quantization and Synthesis
% Methodology". Closes R4-7 ("Add details about ASIC power methodology").
%
% Three factual corrections to the submitted text, all verifiable from the
% scripts in the artifact:
%   S-6  the library is Cadence GPDK045 (gsclib045), NOT TSMC 45 nm
%   S-7  the corner is fast_vdd1v2 @ PVT_1P32V_0C -- 1.32 V, 0 C -- not the
%        claimed "nominal supply voltage of 1.2 V"
%   S-8  power comes from report_power with no SAIF/VCD back-annotation, i.e.
%        default switching statistics, not simulation-derived activity
%
% Stating the corner accurately costs nothing: every design was synthesised
% under the identical setup, so the comparisons stand. Claiming a nominal
% corner we did not use would not survive a reader opening the scripts.

\subsection{Quantisation and Synthesis Methodology}
\label{sec:method}

All datapaths are fixed-point. Operands and intermediate values use a
$Q(1,F)$ format --- one sign bit, one integer bit and $F$ fractional bits,
for a word length $W = F + 2$ --- with $F$ and the stage count $N$ exposed as
RTL parameters so the precision--area surface can be swept directly. The
inputs are block-normalised before the rotator (\S\ref{sec:frontend}), which
is what allows a two-integer-bit format to carry activations whose raw
magnitudes reach $|q| \approx 132$.

Synthesis used Cadence Genus 21.10 targeting the \emph{Cadence GPDK045}
45\,nm standard-cell library (\texttt{gsclib045}), with the high- and
low-$V_t$ variants enabled alongside the nominal-$V_t$ cells. Physical
libraries and the QRC technology file were read in, so areas are
physical-library rather than wire-load-model estimates. Every design ran the
same \texttt{syn\_generic}/\texttt{syn\_map}/\texttt{syn\_opt} flow at medium
effort, under a 2.0\,ns clock constraint (500\,MHz) with 0.2\,ns input and
output delays; the purely combinational blocks are constrained by a matching
2.0\,ns maximum input-to-output delay. Reported areas are total cell area
after \texttt{syn\_opt}, and timing slack is taken from the same report.

Two aspects of the setup are worth stating explicitly, since both bear on how
the absolute numbers should be read.

First, the operating condition is the library's \emph{fast} corner,
\texttt{fast\_vdd1v2} at PVT\_1P32V\_0C --- 1.32\,V and 0\,\textdegree C.
This is optimistic relative to a nominal or worst-case corner, so the absolute
areas and frequencies here are not directly comparable to figures reported at
typical conditions elsewhere in the literature. Every design in this paper,
including the LUT, multiplier and standard-CORDIC baselines, was synthesised
under the identical corner and constraint, so the relative comparisons that
the argument rests on are unaffected.

Second, power figures come from Genus \texttt{report\_power} without SAIF or
VCD back-annotation, and therefore reflect the tool's default switching
statistics rather than activity captured from a RoPE workload. We report them
because they are consistent across designs, but they should be read as a
structural comparison rather than a workload measurement. The direction of
the resulting bias is worth noting: default statistics assume uniform
switching, which understates the advantage of the shift-add datapaths here,
since a multiplier array toggles considerably more than a shifter network
under real data. Annotated power would, if anything, widen the gaps reported
in Table~\ref{tab:perf}.

The RTL, synthesis scripts, constraint files and the evaluation harness are
released at~\cite{github}, including the exact library and corner settings
used, so these choices can be checked rather than taken on trust.
